\section{Related Work}
\label{sec:related_work}

Intrinsic motivation and curiosity provide reinforcement learning agents with internally generated reward signals that guide exploration when extrinsic feedback is sparse or delayed.
Early computational accounts formalize curiosity in terms of progress in prediction or compression, encouraging experience that is learnable but not yet mastered \cite{schmidhuber-1991,schmidhuber-2007,oudeyer-2007}.
In deep RL, intrinsic rewards are often combined with extrinsic reward shaping and trained end-to-end with the policy \cite{ng-1999}.

Exploration has also been studied through count-based novelty and pseudo-counts \cite{bellemare-2016}, optimism and uncertainty estimation \cite{strehl-2006,osband-2016}, and information-theoretic objectives \cite{houthooft-2016}.
These approaches can offer principled exploration incentives but may require accurate density models or calibrated uncertainty estimates in high-dimensional observation spaces.

Prediction-error bonuses are a popular intrinsic signal in high-dimensional RL.
ICM \cite{pathak-2017} uses a learned forward model in a learned feature space, while RND \cite{burda-2018B} uses a learned predictor network against a fixed random target.
Such methods can overexplore in environments with stochastic transitions or observation noise, since irreducible error remains high even when dynamics are not learnable.

Learning-progress methods address this issue by rewarding improvement rather than error magnitude.
Early work measures progress as reduction in prediction error over time \cite{schmidhuber-1991,oudeyer-2007}.
In RL, progress has been localized via region partitioning and bandit-style allocation (e.g., RIAC) \cite{baranes-2009}, encouraging exploration in areas where the agent is currently learning.
The GLPE family builds on this idea by computing region-local learning progress in a learned latent space and combining it with feature-space impact, with an additional gating mechanism to suppress unproductive curiosity when progress stalls.
Section~\ref{sec:method} formalizes these components and Section~\ref{sec:experimental_setup} evaluates them in modern PPO-based training.
